% =====================================================================
%  Kapitel 2 – Ziel der Arbeit (~2 Seiten)
% =====================================================================
\chapter{Ziel der Arbeit}
\label{kap:ziel}

\section{Zielstellung}
\label{sec:zielstellung}

\todoi{Ausgangspunkt Neffy~2.0 (Vorgängerbericht zitieren) und Ziele
von Neffy~3.0 im Gesamtprojekt: sensorbasierte Interaktion, besserer
Touch, Auswahl von Atemübungen per Interaktion.}

\todoi{Ursprünglicher Auftrag Team Sensorik: Anwesenheitserkennung
(Thermalsensor Omron D6T-1A-02, \ac{I2C}) und Atemerkennung (Velostat,
analog am ESP32-S3).}

\todoi{Wendepunkt: Übernahme des kompletten neuen Gehäuses durch das
Team Sensorik -- wann und warum?}

\section{Anforderungsdefinition}
\label{sec:anforderungen}

Tabelle~\ref{tab:anforderungen} fasst die Anforderungen zusammen, gegen
die in Kapitel~\ref{kap:diskussion} evaluiert wird.

\begin{table}[htbp]
  \centering
  \caption{Anforderungen an Neffy~3.0 (Team Sensorik)}
  \label{tab:anforderungen}
  \small
  \begin{tabular}{@{}l p{3.2cm} p{1.6cm} p{6cm}@{}}
    \toprule
    Nr. & Anforderung & Art & Prüfkriterium \\
    \midrule
    A1 & Anwesenheitserkennung & funktional & \todoi{prüfbares Kriterium, z.\,B. Distanz/Reaktionszeit} \\
    A2 & Atemerkennung         & funktional & \todoi{prüfbares Kriterium} \\
    A3 & Touch-Eingabe         & funktional & \todoi{prüfbares Kriterium} \\
    A4 & Gehäuse               & nicht-funktional & \todoi{prüfbares Kriterium, z.\,B. Bauraum, Montage} \\
    \bottomrule
  \end{tabular}
\end{table}
\todo{Anforderungsliste ist ein Entwurf aus dem bekannten Projektstand -- bitte abgleichen und ergänzen}

\section{Übergeordnete Fragestellung}
\label{sec:fragestellung}

\begin{center}
  \fbox{\parbox{0.85\linewidth}{\centering\itshape
    \todoi{Übergeordnete Fragestellung als hervorgehobene Frage formulieren.}}}
\end{center}

\section{Projektorganisation und Aufgabenverteilung}
\label{sec:organisation}

\todoi{Kurz: drei Teams (Sensorik, Programmierung, Design), Arbeitsweise,
Abstimmung, Termine.}

\begin{table}[htbp]
  \centering
  \caption{Aufgabenverteilung im Team Sensorik}
  \label{tab:aufgaben}
  \small
  \begin{tabular}{@{}p{5cm} p{4.5cm} p{4.5cm}@{}}
    \toprule
    Aufgabe & Badre & \todoi{Name Teampartner:in} \\
    \midrule
    Anwesenheitserkennung (Thermalsensor) & \todoi{Anteil} & \todoi{Anteil} \\
    Atemerkennung (Velostat)              & \todoi{Anteil} & \todoi{Anteil} \\
    Auswahl und Test Drucksensoren        & \todoi{Anteil} & \todoi{Anteil} \\
    Gehäuse: \ac{CAD}-Konstruktion        & \todoi{Anteil} & \todoi{Anteil} \\
    Gehäuse: 3D-Druck und Nachbearbeitung & \todoi{Anteil} & \todoi{Anteil} \\
    Software Signalverarbeitung           & \todoi{Anteil} & \todoi{Anteil} \\
    \bottomrule
  \end{tabular}
\end{table}

\todoi{Schnittstellen zu Team Programmierung und Team Design
(nur soweit für Sensorik/Gehäuse relevant).}
